
\documentclass[10pt,letterpaper,unboxed,cm]{article}
\usepackage[margin=1in]{geometry}
\usepackage{graphicx}
\usepackage{enumerate, comment}
\usepackage{adjustbox}
\usepackage{amsmath,mathabx}

\newcommand{\st}{~\mid~}
\newcommand{\ind}{$~~~~~$}




\begin{document}

\hfill{CSE 21 Fall 2026}

\hfill{Homework 3}

\hfill{Due date: Monday October 19, 2026 at 11:59pm}

\begin{center}
\begin{minipage}[t]{5.7in}
\rule{\linewidth}{2pt}
{\sc Instructions}\newline

This homework should be done in {\bf groups} of one to four students, without assistance from anyone besides the instructional staff and your group members.  Homework must be submitted through Gradescope by a {\bf single representative} of your group and received by {\bf 11:59pm} on the due date. There are no exceptions to this rule.\\
 
You will be able to look at your scanned work before submitting it. 
 You must {\bf type} your solutions. (hand-drawn diagrams are okay.)
 Your group representative can resubmit your assignment as many times as you like before the deadline. Only the most recent submission will be graded.\\

Students should consult their textbook, class notes, lecture slides, podcasts, group members, instructors, TAs, and tutors when 
they need help with homework. You may ask questions about the homework in office hours, but questions on Piazza should be private, visible only to instructors. 

This assignment will be graded for not only the \emph{correctness} of your answers, but on your ability to present your ideas clearly and logically. You should explain or justify, present clearly how you arrived at your conclusions and justify the correctness of your answers with mathematically sound reasoning (unless explicitly told not to). Whether you use formal proof techniques or write a more informal argument for why something is true, your answers should always be well-supported. Your goal should be to {\bf convince the reader} that your results and methods are sound.






{\sc Key Concepts}  basic probability principles,
independence, conditional probability, Bayes’ Theorem, Random Sampling, Random Variables, .



\rule{\linewidth}{2pt}
\end{minipage} \hfill

\end{center}

\begin{enumerate}


\item (20 points)
Consider the set $H$ of all non-negative integer solutions to the equation:

$$a_1 + a_2 + a_3 + a_4 + a_5 + a_6 = 30$$

(You can think of $H$ as a set of 6-tuples of non-negative integers that add up to 30).

You wish to pick one of the elements of $H$ at random and the way you decide to do it is the following:

Initialize a 6-tuple to be all zeros (0,0,0,0,0,0).

Roll a fair 6 sided die 30 times. For each roll that shows a 1, increment the 1st position, for each roll that shows a 2, increment the 2nd position and so on. (include a justification of each problem.)

\begin{enumerate}
\item
What is the probability that you get the element: $(5,5,5,5,5,5)$?
\item
What is the probability that at least one of the entries is $0$?
\item
What is the probability that one of the entries is exactly 16?
\item
Explain why this is not a uniform random sampling. (Show that the probability of each outcome is not equally likely.)
\item
Which elements are the least likely to occur and why?
\end{enumerate}




\item  (12 points)
Suppose you are at a casino and there are three slot machines. You know that two slot machines each win with probability 0.5, and the other wins with a probability of 0.9.

You pick one of the slot machines (equally at random) and you play 8 games and win them all.

\begin{itemize}
\item
What is the probability that you are at the slot machine that wins with probability of 0.9?

\item
If you decide to stay at the same slot machine (after winning 8 games), what is the probability that your next game results in a win?
\end{itemize}

\item
\emph{Partnership Dominos} is a game that involves 4 players and a \emph{Double-6 Domino set}.

A \emph{Double-6 Domino set} is a set of 28 different tiles such that each tile has two values (one on each side. The values are given by a number of \emph{pips} or dots to signify the value.) There is one domino for each combination of values from 0 through 6 (including doubles). Here is a picture of a full set of 28 Double-6 dominos:

\includegraphics[scale=0.4]{dominos}

At the start of the round, each of the 4 players (North, South, East, West) gets 7 tiles.

\begin{enumerate}

\item
What is the probability that North has the double-six tile (the tile with 6 pips on either side)?


\item
What is the probability that North has at least one double? (there are 7 doubles)

\item
What is the probability that North has at least one domino with a 6 and at least one domino with a blank?

\item
What is the probability that North and South together hold all 7 tiles containing the value 6?

\item
A ``Bad hand" is a set of 7 dominos with at least four doubles. What is the probability that one of the players has a bad hand?



\end{enumerate}

\item
\emph{Fantasy Five} is a lottery game you can play for \$1. The way it works is, from the set {1,2,3,...,39}, you select a subset of 5 different numbers.

The lottery commissioner also selects 5 \emph{winning numbers} uniformly at random.

You win money based on how many of your numbers match with the \emph{winning numbers}.

Here are the payouts:

5 out of 5 pays \$67000\\
4 out of 5 pays \$401\\
3 out of 5 pays \$15\\
2 out of 5 pays \$1\\
1 out of 5 pays \$0\\
0 out of 5 pays \$0\\


\begin{enumerate}
\item (4 points)
What is the expected number of matches you make with the winning numbers?

\item (4 points)
What is the expected profit you make from playing this game?

\item (4 points)
What is the expected number of times you must play this game before and including your first win of at least one dollar?
\end{enumerate}



\end{enumerate}

\end{document}
